This was prompted by observing that both cases are based on the same index-theoretic technique, namely an index theorem for certain \emph{Callias-type} operators, suggesting that there should be a more general geometric statement underlying both.
This is achieved in Theorem~\ref{thm:quantitative-codim1-obstr} below.
We show that in the presence of a~hypersurface \(M \subset X\) with an index-theoretic obstruction, positive scalar curvature \(\geq \sigma > 0\) in a~neighborhood \(\{p \in X \mid d(p, M) \leq d/2 \}\) of width \(d\) forces the scalar curvature somewhere else to be negative in a~quantifiable way depending on \(d\) and \(\sigma\).
We formulate this for the normalized case \(\sigma = n(n-1) = \scal_{\Sphere^n}\).

\begin{Theorem} \label{thm:quantitative-codim1-obstr}
 Let \(X\) be a spin \(n\)-manifold without boundary and \(M \subseteq X\) a closed codimension one submanifold with trivial normal bundle such that the induced map \(\pi_1 M \to \pi_1 X\) is injective.
 Suppose that the Rosenberg index \(\alpha(M) \in \KO_{n-1}(\Cstar \pi_1 M)\) does not vanish.

 Let \(g\) be a complete Riemannian metric on \(X\) which has scalar curvature bounded below by \(n(n-1)\) in a neighborhood of \(M\) of width \(d < 2\pi/n\).
 Then
 \begin{equation} \label{eq:negativity}
 \inf_{p \in X} \scal_g(p) \leq - n(n-1) \tan\left(\frac{nd}{4}\right)^2.
 \end{equation}
\end{Theorem}
